\cotefile{31}
\scanpage{70}
\uline{Corollaire} Supposons que pour \unsure{tt.} \deleted{\unsure{partie}} \inserted{sous-schéma} fermé
\deleted{\unsure{irréductible}} $Y$ de $X$, et tt \deleted{$y\in Y$} \unsure{pt} \uline{\unsure{géom}}
$\bar y$ de $Y$, \unsure{les|l'} \unsure{anneau|anneaux} total des fractions $K_{Y,\bar y}$ de
$\underline O_{Y,\bar y}$ \deleted{\unsure{soit de|soient de} $\operatorname{cd}_{\ell}$} \unsure{satisfasse|satisfont} $\operatorname{cd}_{\ell}K_{Y,\bar y}\leq\dim\underline O_{Y,\bar y}$
$[=\dim\underline O_{Y,y}]$. \deleted{\uline{Alors} \unsure{soit} $\varphi$ \unsure{définie} \unsure{en}} [N.B. C'est \unsure{probablement}
tjrs. le cas, du moins si $X$ est ``bon''].

Posant pour \unsure{tt} \deleted{\unsure{e} $x\in X$ \unsure{partie} $Y$} $x\in X$
\[
\boxed{\varphi(x)=\sup_{y\in\bar x}\bigl(\operatorname{cd}_{\ell}(k(y))+2\dim\underline O_{\bar x,y}\bigr)}
\]
\deleted{Alors} et pour tte partie \inserted{fermée} $Y$ de $X$
\[
\deleted{\operatorname{cd}_{\ell}(X)\leq\unsure{?}}
\]
\[
\varphi(Y)=\sup_{y\in Y}\varphi(y)=\sup_{y\in Y}\bigl(\operatorname{cd}_{\ell}(k(y))+2\dim\underline O_{Y,y}\bigr)
\]
Alors
\[
\operatorname{cd}_{\ell}(X)\leq\varphi(X)
\]
% slanted
\margin{$\unsure{\psi|\varphi}(x)=\operatorname{cd}_{\ell}(x)+\deleted{\unsure{\sup_{y\in\bar x}}}\dim\unsure{\bar x|\bar X}_{\deleted{\unsure{y}}}$}

% slanted
\margin{\uline{Question} $A$ local noeth. hens. intègre, \unsure{de corps} résiduel $k$ \deleted{de dim \unsure{finie}}, tel que $\operatorname{cd}_{\ell}(k)<+\infty$. Alors $\operatorname{cd}_{\ell}K\geq\dim A+\operatorname{cd}_{\ell}k$\unsure{?} ??}

\uline{Exemples} 1) Schémas algébriques, plus \unsure{généralement|générlmt}
extensions résiduelles de degré de transcendance
\unsure{finies|finis|fini} [NB. ce degré de transcendance varie
simplement avec la dimension \ldots]

2) Corps de nombres

3) (conjecturalement) Ouvert de $\operatorname{Spec} A$,
$A$ \unsure{strictement} local

4) Schéma \unsure{alg.|rég.} \uline{affine} de dim 1
% archive sheet number: 69
